\uuid{Fnys}
\titre{Énoncé du théorème central limite}
\auteur{Maxime Nguyen}
\organisation{AMSCC}
\datecreate{2026-11-27}
\niveau{L2}
\module{Probabilité}
\chapitre{Probabilité continue}
\sousChapitre{Théorème central limite}
\difficulte{2}
\theme{théorème central limite, convergence en loi, loi normale}

\contenu{
\begin{question}
On note $\mathcal{N}(\mu,\sigma)$ la loi normale d'espérance $\mu$ et d'écart-type $\sigma$. Soit $(X_k)_{k \geq 1}$ une suite de variables indépendantes de même loi, d'espérance $\mu$ et de variance $\sigma^2 > 0$, et $S_n = X_1 + \dots + X_n$. Quelles affirmations sont vraies ?
\begin{answers}
\bad{Pour $n$ grand, $S_n$ suit approximativement la loi $\mathcal{N}(n\mu, n\sigma)$.}
\good{$\dfrac{S_n - n\mu}{\sigma\sqrt n}$ converge en loi vers une variable de loi $\mathcal{N}(0,1)$.}
\bad{Ce résultat n'est valable que si les $X_k$ suivent une loi normale.}
\good{Pour $n$ grand, $S_n$ suit approximativement la loi $\mathcal{N}(n\mu, \sigma\sqrt n)$.}
\good{Si les $X_k$ suivent la loi $\mathcal{N}(\mu,\sigma)$, alors $\dfrac{S_n - n\mu}{\sigma\sqrt n}$ suit exactement la loi $\mathcal{N}(0,1)$ pour tout $n$.}
\end{answers}
\begin{explanations}
B est l'énoncé du cours, D sa lecture pratique : $\mathbb{E}(S_n) = n\mu$ et $\mathbb{V}(S_n) = n\sigma^2$, donc l'écart-type de $S_n$ est $\sigma\sqrt n$ (et A est faux).

C est faux : c'est tout l'intérêt du théorème, il s'applique à n'importe quelle loi ayant une variance (Bernoulli, exponentielle, uniforme…).

E : dans le cas normal, $S_n \sim \mathcal{N}(n\mu, \sigma\sqrt n)$ exactement, par la règle de la somme ; il n'y a rien à approcher.

Contre-sens visés : multiplier l'écart-type par $n$ au lieu de $\sqrt n$
(A) ; croire que le théorème suppose des variables normales (C).
\end{explanations}
\end{question}
}
